\paragrafos{Μήκος ευθύγραμμου τμήματος}{Β.1.2, «Μήκος ευθύγραμμου τμήματος»}

\begin{anagke}
Για να σχεδιάσουν τον κήπο, η παρέα μετράει πρώτα το οικόπεδο. Ο Στέλιος μετράει το μήκος του μονοπατιού με βήματα και βρίσκει $14$. Ο Ηλίας το μετράει με τα δικά του βήματα και βρίσκει $16$. Ποιος έχει δίκιο; «Και οι δύο», λέει η Πολυξένη: «το βήμα του Στέλιου είναι μεγαλύτερο. Έχουμε διαφορετικές μονάδες». Η Αγγελική παίρνει τη μετροταινία του παππού: $9{,}8$ μέτρα. «Αυτό είναι σίγουρο: όλοι θα βρούμε το ίδιο νούμερο».

Μετά η Ρήνη ρωτάει: «Θέλουμε το συντριβάνι να είναι ακριβώς στη μέση του μονοπατιού. Πώς βρίσκουμε τη μέση;». Ο Νίκος πιάνει ένα σχοινί με το μήκος του μονοπατιού και το διπλώνει στα δύο: «Εκεί που διπλώνει είναι η μέση!». Ο Βασίλης προσθέτει: «Και για τον φράχτη γύρω από το παρτέρι πρέπει να προσθέσουμε όλες τις πλευρές του». Η Φανή τα σημειώνει όλα στο τετράδιό της και ρωτάει: «Μα πόσα μέτρα φράχτη είναι αυτά, αν κάποιες πλευρές τις μετρήσαμε σε εκατοστά;».

Σε αυτή την παράγραφο μαθαίνουμε να μετράμε το μήκος ενός τμήματος με σωστές μονάδες, να μετατρέπουμε μονάδες, να βρίσκουμε το μέσο ενός τμήματος και να προσθέτουμε και να αφαιρούμε τμήματα. Έτσι θα υπολογίζουμε και την περίμετρο ενός σχήματος.
\end{anagke}

\begin{lexi}{απόσταση, μέσο, περίμετρος}{απόσταση, μέσο, περίμετρος}
Η \lexiword{απόσταση} δύο σημείων είναι το πόσο «απέχουν». Το \lexiword{μέσο} ενός τμήματος είναι το σημείο που βρίσκεται «στη μέση». Η \lexiword{περίμετρος} είναι το «γύρω μέτρημα» (περί + μέτρο): το συνολικό μήκος του περιγράμματος ενός σχήματος. Το \lexiword{μέτρο} στα αρχαία ελληνικά είναι το «μέτρον», η ποσότητα με την οποία μετράμε.
\end{lexi}

\begin{metafrasi}
Το μήκος του τμήματος $\mathrm{ΑΒ}$ είναι $6{,}4$ εκατοστά. & $\mathrm{ΑΒ}=6{,}4\,\text{cm}$\\
Τα δύο τμήματα έχουν το ίδιο μήκος. & $\mathrm{ΑΒ}=\Gamma\Delta$\\
Η απόσταση των $\mathrm{Α}$ και $\mathrm{Β}$. & το μήκος του $\mathrm{ΑΒ}$\\
Το $\mathrm{Μ}$ είναι το μέσο του $\mathrm{ΑΒ}$. & $\mathrm{ΑΜ}=\mathrm{ΜΒ}$\\
Το $\mathrm{Β}$ είναι ανάμεσα στα $\mathrm{Α}$ και $\Gamma$. & $\mathrm{ΑΓ}=\mathrm{ΑΒ}+\mathrm{ΒΓ}$\\
Το άθροισμα των πλευρών. & περίμετρος\\
\end{metafrasi}

\tmima{Η θεωρία}

\tmimasmall{Μέτρηση μήκους}

Για να συγκρίνουμε δύο μεγέθη, δηλαδή για να δούμε αν το ένα είναι μεγαλύτερο, μικρότερο ή ίσο με το άλλο, τα συγκρίνουμε με ένα σταθερό μέγεθος που λέγεται \textbf{μονάδα μέτρησης}.

\begin{kanonas}{μήκος ευθύγραμμου τμήματος}
Βρίσκουμε το \textbf{μήκος} ενός ευθύγραμμου τμήματος συγκρίνοντάς το με ένα άλλο ευθύγραμμο τμήμα που το λέμε \textbf{μονάδα μήκους}. Η βασική μονάδα μήκους που χρησιμοποιούμε είναι το \textbf{μέτρο} (m).
\end{kanonas}

Για μικρά μήκη χρησιμοποιούμε συνήθως το \textbf{υποδεκάμετρο}, που είναι το ένα δέκατο του μέτρου. Αν το τμήμα είναι πολύ μεγαλύτερο ή πολύ μικρότερο από το μέτρο, διαλέγουμε ένα πολλαπλάσιο ή μια υποδιαίρεση του μέτρου.

\tmimasmall{Μονάδες μέτρησης μήκους}

\renewcommand{\arraystretch}{1.25}
\begin{center}
\begin{tabular}{@{}lll@{}}
\toprule
 & \textbf{Ονομασία} & \textbf{Σύμβολο}\\
\midrule
Πολλαπλάσιο του μέτρου & Χιλιόμετρο & km\\
Βασική μονάδα & Μέτρο & m\\
Υποδιαιρέσεις του μέτρου & Δεκατόμετρο & dm\\
 & Εκατοστόμετρο & cm\\
 & Χιλιοστόμετρο & mm\\
\bottomrule
\end{tabular}
\end{center}

\begin{center}
\begin{tikzpicture}[x=1cm,y=1cm]
\foreach \x/\n in {0/km,2.6/m,5.2/dm,7.8/cm,10.4/mm}{\node[draw,rounded corners=3pt,minimum width=1.5cm,fill=gblueL] (\n) at (\x,0) {\n};}
\draw[-{Stealth},gblue] (km.north) .. controls +(0.4,0.6) and +(-0.4,0.6) .. node[above] {\small $\times1000$} (m.north);
\draw[-{Stealth},gblue] (m.north) .. controls +(0.4,0.6) and +(-0.4,0.6) .. node[above] {\small $\times10$} (dm.north);
\draw[-{Stealth},gblue] (dm.north) .. controls +(0.4,0.6) and +(-0.4,0.6) .. node[above] {\small $\times10$} (cm.north);
\draw[-{Stealth},gblue] (cm.north) .. controls +(0.4,0.6) and +(-0.4,0.6) .. node[above] {\small $\times10$} (mm.north);
\draw[-{Stealth},gwarmD] (mm.south) .. controls +(-0.4,-0.6) and +(0.4,-0.6) .. node[below] {\small $:10$} (cm.south);
\draw[-{Stealth},gwarmD] (cm.south) .. controls +(-0.4,-0.6) and +(0.4,-0.6) .. node[below] {\small $:10$} (dm.south);
\draw[-{Stealth},gwarmD] (dm.south) .. controls +(-0.4,-0.6) and +(0.4,-0.6) .. node[below] {\small $:10$} (m.south);
\draw[-{Stealth},gwarmD] (m.south) .. controls +(-0.4,-0.6) and +(0.4,-0.6) .. node[below] {\small $:1000$} (km.south);
\end{tikzpicture}
\end{center}

Από τη μεγαλύτερη προς τη μικρότερη μονάδα \textbf{πολλαπλασιάζουμε}: $1\,\text{m}=10\,\text{dm}=100\,\text{cm}=1000\,\text{mm}$ και $1\,\text{km}=1000\,\text{m}$. Από τη μικρότερη προς τη μεγαλύτερη \textbf{διαιρούμε}: $1\,\text{mm}=0{,}1\,\text{cm}=0{,}001\,\text{m}$ και $1\,\text{m}=0{,}001\,\text{km}$.

Στην Αγγλία, την Αμερική και σε μερικές ακόμη χώρες η βασική μονάδα μήκους είναι η \textbf{υάρδα} ή γιάρδα (yd). Η $1$ yd διαιρείται σε $3$ πόδια (ft) και ισούται περίπου με $0{,}9144$ m. Το $1$ πόδι διαιρείται σε $12$ ίντσες (in) και ισούται περίπου με $0{,}3048$ m.

\tmimasmall{Ίσα τμήματα, απόσταση, μέσο}

Δύο τμήματα είναι \textbf{ίσα} όταν έχουν το ίδιο μήκος, όταν δηλαδή το ένα εφαρμόζει ακριβώς στο άλλο (π.χ. με τον διαβήτη). Για να συγκρίνουμε δύο μήκη πρέπει να τα εκφράσουμε στην ίδια μονάδα.

\begin{orismos}{απόσταση δύο σημείων}
\textbf{Απόσταση} δύο σημείων $\mathrm{Α}$ και $\mathrm{Β}$ λέγεται το μήκος του ευθύγραμμου τμήματος $\mathrm{ΑΒ}$ που τα ενώνει.
\end{orismos}

\begin{orismos}{μέσο ευθύγραμμου τμήματος}
\textbf{Μέσο} ενός ευθύγραμμου τμήματος $\mathrm{ΑΒ}$ ονομάζουμε το σημείο $\mathrm{Μ}$ του τμήματος που απέχει εξίσου από τα άκρα του: $\mathrm{ΑΜ}=\mathrm{ΜΒ}$.
\end{orismos}

\begin{center}
\begin{tikzpicture}
\draw[thick] (0,0)--(6,0);
\foreach \x/\l in {0/Α,3/Μ,6/Β}{\fill (\x,0) circle (2pt); \node[below] at (\x,-0.1) {$\mathrm{\l}$};}
\node[gblue,above] at (1.5,0.05) {\small $2$ cm}; \node[gblue,above] at (4.5,0.05) {\small $2$ cm};
\end{tikzpicture}
\end{center}

Κάθε τμήμα έχει ένα μοναδικό μέσο, και $\mathrm{ΑΜ}=\mathrm{ΜΒ}=\dfrac{\mathrm{ΑΒ}}{2}$.

\tmimasmall{Άθροισμα και διαφορά τμημάτων}

Αν πάρουμε σε μια ευθεία με τη σειρά τα σημεία $\mathrm{Α}$, $\mathrm{Β}$ και $\Gamma$, τότε:
\begin{itemize}[leftmargin=1.2em,itemsep=0.3ex,topsep=2pt]
\item το τμήμα $\mathrm{ΑΓ}$ είναι το \textbf{άθροισμα} των $\mathrm{ΑΒ}$ και $\mathrm{ΒΓ}$: $\mathrm{ΑΓ}=\mathrm{ΑΒ}+\mathrm{ΒΓ}$.
\item το τμήμα $\mathrm{ΑΒ}$ είναι η \textbf{διαφορά} των $\mathrm{ΑΓ}$ και $\mathrm{ΒΓ}$: $\mathrm{ΑΒ}=\mathrm{ΑΓ}-\mathrm{ΒΓ}$.
\end{itemize}

\begin{center}
\begin{tikzpicture}
\draw[thick] (0,0)--(7,0);
\foreach \x/\l in {0/Α,2.5/Β,7/Γ}{\fill (\x,0) circle (2pt); \node[below] at (\x,-0.1) {$\mathrm{\l}$};}
\draw[gblue,thick] (0,0.4)--(2.5,0.4); \draw[gwarmD,thick] (2.5,0.4)--(7,0.4);
\node[gblue,above] at (1.25,0.45) {$\mathrm{ΑΒ}$}; \node[gwarmD,above] at (4.75,0.45) {$\mathrm{ΒΓ}$};
\draw[thick,ggray] (0,-0.9)--(7,-0.9); \node[ggray,below] at (3.5,-0.95) {$\mathrm{ΑΓ}=\mathrm{ΑΒ}+\mathrm{ΒΓ}$};
\end{tikzpicture}
\end{center}

Για να κατασκευάσουμε με τον διαβήτη ένα τμήμα με μήκος $\mathrm{ΑΒ}+\Gamma\Delta$, σχεδιάζουμε μια ευθεία, παίρνουμε πάνω της ένα τμήμα ίσο με το $\mathrm{ΑΒ}$ και στη συνέχεια, από το άκρο του, ένα τμήμα ίσο με το $\Gamma\Delta$. Για τη διαφορά $\mathrm{ΑΒ}-\Gamma\Delta$ (όταν $\mathrm{ΑΒ}>\Gamma\Delta$) παίρνουμε πάνω στο $\mathrm{ΑΒ}$, από το ένα άκρο του, τμήμα ίσο με το $\Gamma\Delta$ και το κομμάτι που μένει είναι η διαφορά.

\tmimasmall{Περίμετρος ευθύγραμμου σχήματος}

\textbf{Κυρτό} ονομάζεται ένα ευθύγραμμο σχήμα όταν η προέκταση κάθε πλευράς του αφήνει το σχήμα στο ίδιο ημιεπίπεδο. Διαφορετικά ονομάζεται \textbf{μη κυρτό}.

\begin{center}
\begin{tikzpicture}[scale=0.9]
\coordinate (A) at (0,0); \coordinate (B) at (2.5,0); \coordinate (C) at (3.1,1.8); \coordinate (D) at (1.2,2.8); \coordinate (E) at (-0.5,1.5);
\draw[thick,fill=gblueL] (A)--(B)--(C)--(D)--(E)--cycle;
\draw[dashed,ggray] ($(A)!-0.5!(B)$)--($(A)!1.7!(B)$);
\node at (1.3,-0.8) {\small κυρτό};
\begin{scope}[xshift=6cm]
\coordinate (P) at (0,0); \coordinate (Q) at (3,0); \coordinate (R) at (3,2.2); \coordinate (S) at (1.6,0.9); \coordinate (T) at (0,2.2);
\draw[thick,fill=gwarm] (P)--(Q)--(R)--(S)--(T)--cycle;
\draw[dashed,ggray] ($(S)!-0.7!(R)$)--($(S)!2.2!(R)$);
\node at (1.5,-0.8) {\small μη κυρτό};
\end{scope}
\end{tikzpicture}
\end{center}

\begin{orismos}{περίμετρος}
Το άθροισμα των πλευρών ενός ευθύγραμμου σχήματος λέγεται \textbf{περίμετρος} του σχήματος.
\end{orismos}

Για παράδειγμα, αν ένα τετράπλευρο έχει πλευρές $4$ cm, $2{,}6$ cm, $3{,}1$ cm και $3$ cm, η περίμετρός του είναι $4+2{,}6+3{,}1+3=12{,}7$ cm.

\begin{prosoxi}
\begin{itemize}[leftmargin=1.2em,itemsep=0.3ex,topsep=2pt]
\item Για να συγκρίνω ή να προσθέσω μήκη πρέπει πρώτα να τα φέρω στην ίδια μονάδα: τα $0{,}8$ m δεν είναι μικρότερα από τα $75$ cm, γιατί $0{,}8$ m $=80$ cm.
\item Στις μετατροπές μεγάλη $\to$ μικρή μονάδα πολλαπλασιάζω, και μικρή $\to$ μεγάλη διαιρώ. Αν ο αριθμός που βγάζω είναι μεγαλύτερος για τη μεγαλύτερη μονάδα, κάτι πήγε στραβά.
\item Το μέσο \emph{ανήκει} στο τμήμα: δεν αρκεί να απέχει εξίσου από τα άκρα.
\item Η σχέση $\mathrm{ΑΓ}=\mathrm{ΑΒ}+\mathrm{ΒΓ}$ ισχύει όταν το $\mathrm{Β}$ βρίσκεται \emph{ανάμεσα} στα $\mathrm{Α}$ και $\Gamma$ πάνω στην ίδια ευθεία.
\item Η περίμετρος μετράει μόνο το περίγραμμα, όχι την επιφάνεια του σχήματος.
\end{itemize}
\end{prosoxi}

\tmima{Λυμένα παραδείγματα}

\paradeigma{Μετατροπές μονάδων}
\tixero $2{,}4$ km, $35$ cm και $7$ dm.\par
\tithelo Να τα μετατρέψω σε m, m και mm αντίστοιχα.\par
\possindeo Από km σε m πολλαπλασιάζω επί $1000$. Από cm σε m διαιρώ με $100$. Από dm σε mm πολλαπλασιάζω επί $100$.\par
\apantisi $2{,}4$ km $=2400$ m, $35$ cm $=0{,}35$ m, $7$ dm $=700$ mm.\par
\elegxos Το $0{,}35$ m είναι λιγότερο από ένα μέτρο, όπως τα $35$ cm.

\paradeigma{Ποιο είναι μεγαλύτερο;}
\tixero Δύο ράφια έχουν μήκη $0{,}8$ m και $75$ cm.\par
\tithelo Ποιο είναι μεγαλύτερο;\par
\possindeo Τα φέρνω στην ίδια μονάδα: $0{,}8$ m $=80$ cm.\par
\apantisi Μεγαλύτερο είναι το πρώτο ($80$ cm $>75$ cm).\par
\elegxos Σε mm: $800>750$.

\paradeigma{Το συντριβάνι}
\tixero Το μονοπάτι είναι $9{,}8$ μ. και το συντριβάνι θα μπει στο μέσο του.\par
\tithelo Πόσο απέχει το συντριβάνι από κάθε άκρο.\par
\possindeo Το μέσο $\mathrm{Μ}$ έχει $\mathrm{ΑΜ}=\dfrac{\mathrm{ΑΒ}}{2}=\dfrac{9{,}8}{2}$.\par
\apantisi $4{,}9$ μ. από κάθε άκρο.\par
\elegxos $4{,}9+4{,}9=9{,}8$.

\paradeigma{Τα δέντρα στη σειρά}
\tixero Τα δέντρα $\mathrm{Α},\,\mathrm{Β},\,\Gamma,\,\Delta$ είναι με αυτή τη σειρά σε ευθεία και $\mathrm{ΑΒ}=3$ μ., $\mathrm{ΒΓ}=4$ μ., $\Gamma\Delta=2{,}5$ μ.\par
\tithelo Τα $\mathrm{ΑΓ}$, $\mathrm{ΑΔ}$ και $\mathrm{ΒΔ}$.\par
\possindeo $\mathrm{ΑΓ}=\mathrm{ΑΒ}+\mathrm{ΒΓ}=3+4=7$. $\mathrm{ΑΔ}=\mathrm{ΑΓ}+\Gamma\Delta=7+2{,}5=9{,}5$. $\mathrm{ΒΔ}=\mathrm{ΑΔ}-\mathrm{ΑΒ}=9{,}5-3=6{,}5$.\par
\apantisi $\mathrm{ΑΓ}=7$ μ., $\mathrm{ΑΔ}=9{,}5$ μ. και $\mathrm{ΒΔ}=6{,}5$ μ.\par
\elegxos $\mathrm{ΒΔ}=\mathrm{ΒΓ}+\Gamma\Delta=4+2{,}5=6{,}5$.

\paradeigma{Περίμετρος με μικτές μονάδες}
\tixero Ένα τρίγωνο έχει πλευρές $4$ cm, $35$ mm και $0{,}06$ m.\par
\tithelo Η περίμετρός του σε cm.\par
\possindeo Τα φέρνω όλα σε cm: $35$ mm $=3{,}5$ cm και $0{,}06$ m $=6$ cm. Προσθέτω: $4+3{,}5+6$.\par
\apantisi $13{,}5$ cm.\par
\elegxos Αν είχα προσθέσει $4+35+0{,}06$ χωρίς μετατροπή, θα έβγαινε παράλογος αριθμός.

\paradeigma{Η περίμετρος του παρτεριού}
\tixero Ορθογώνιο παρτέρι με πλευρές $4{,}5$ μ. και $3{,}2$ μ.\par
\tithelo Πόσος φράχτης χρειάζεται για να το περιβάλει, και τι γίνεται αν διπλασιάσουμε κάθε πλευρά.\par
\possindeo Η περίμετρος είναι το άθροισμα των πλευρών: $4{,}5+3{,}2+4{,}5+3{,}2=15{,}4$. Με διπλάσιες πλευρές: $9+6{,}4+9+6{,}4=30{,}8$.\par
\apantisi $15{,}4$ μ. και, με διπλάσιες πλευρές, $30{,}8$ μ., δηλαδή διπλάσια περίμετρος.\par
\elegxos $2\cdot15{,}4=30{,}8$.

\tmima{Ασκήσεις}

\zoni{Το κατάλαβα}{μονάδες, μέσο, άθροισμα, περίμετρος}

\ex Ποια μονάδα (km, m, dm, cm, mm) είναι προτιμότερη για τις μετρήσεις; α) απόσταση δύο πόλεων \quad β) μήκος κρεβατιού \quad γ) πλάτος παραθύρου \quad δ) πάχος τετραδίου \quad ε) μήκος ποδηλατόδρομου.

\ex Μετάτρεψε σε m: α) $1{,}78$ km \quad β) $0{,}12$ dm \quad γ) $21{,}3$ cm \quad δ) $0{,}03$ km \quad ε) $2$ dm \quad στ) $10$ mm.

\ex Μετάτρεψε σε cm: α) $0{,}7$ m \quad β) $45$ mm \quad γ) $1{,}2$ dm \quad δ) $0{,}03$ km.

\ex Σύγκρινε και βάλε σε αύξουσα σειρά τα μήκη: $0{,}35$ m, $3{,}4$ dm, $36$ cm, $345$ mm.

\ex Σχεδίασε ένα ευθύγραμμο τμήμα $\mathrm{ΑΒ}=6{,}5$ cm και βρες το μέσο του $\mathrm{Μ}$. Μέτρησε τα $\mathrm{ΑΜ}$ και $\mathrm{ΜΒ}$ και σύγκρινέ τα.

\ex Τα $\mathrm{Α},\,\mathrm{Β},\,\Gamma$ είναι με αυτή τη σειρά σε μια ευθεία. α) Αν $\mathrm{ΑΒ}=4{,}2$ cm και $\mathrm{ΒΓ}=3{,}5$ cm, βρες το $\mathrm{ΑΓ}$. β) Αν $\mathrm{ΑΓ}=10$ cm και $\mathrm{ΑΒ}=6{,}5$ cm, βρες το $\mathrm{ΒΓ}$.

\ex Σχεδίασε τμήματα $\mathrm{ΑΒ}=5$ cm, $\Gamma\Delta=3$ cm και $\mathrm{ΕΖ}=2$ cm. Κατασκεύασε με τον διαβήτη τμήμα με μήκος: α) $\mathrm{ΑΒ}+\Gamma\Delta$ \quad β) $\mathrm{ΑΒ}-\mathrm{ΕΖ}$ \quad γ) $\mathrm{ΑΒ}+\Gamma\Delta-\mathrm{ΕΖ}$. Έλεγξε με τον χάρακα.

\ex Βρες την περίμετρο: α) ορθογωνίου με πλευρές $6$ και $2{,}5$ \quad β) τριγώνου με πλευρές $5$, $6$, $7{,}5$ \quad γ) τετραγώνου με πλευρά $2{,}8$ \quad δ) πενταγώνου με όλες τις πλευρές ίσες με $4$.

\zoni{Το χρησιμοποιώ}{βήματα, σχοινιά, φράχτες, δρόμοι}

\ex \textbf{Τα βήματα.} Ο Στέλιος μετράει το μονοπάτι με $14$ βήματα και ο Ηλίας με $16$. Γιατί βγαίνουν διαφορετικοί αριθμοί; Αν το μονοπάτι είναι $9{,}8$ μ., πόσο είναι το βήμα του καθενός;

\ex \textbf{Το συντριβάνι.} Ένα άλλο μονοπάτι είναι $12{,}6$ μ. και το συντριβάνι μπαίνει στο μέσο του. Πόσο απέχει από κάθε άκρο; Ένα παγκάκι βρίσκεται πάνω στο μονοπάτι, $1{,}5$ μ. από το συντριβάνι προς το ένα άκρο. Πόσο απέχει το παγκάκι από αυτό το άκρο;

\ex \textbf{Το κορδόνι.} Η θεία Μαρία πουλά κορδόνι για κήπο προς $1{,}80$ € το μέτρο. Η Ρήνη αγοράζει $3$ μ. και $45$ cm. Πόσα μέτρα αγόρασε και πόσα ευρώ πλήρωσε;

\ex \textbf{Ο φράχτης.} Το ορθογώνιο παρτέρι έχει πλευρές $4{,}5$ μ. και $3{,}2$ μ. Ο φράχτης κοστίζει $3$ € το μέτρο και ένα μέτρο μένει ανοιχτό για την είσοδο. Πόσο θα κοστίσει;

\ex \textbf{Η διαδρομή.} Ο Ηλίας πηγαίνει από το σπίτι στο σχολείο με στάσεις: $250$ μ., $180$ μ. και $320$ μ. Πόσα μέτρα περπάτησε και πόσα χιλιόμετρα είναι;

\ex \textbf{Αμερικανικά μέτρα.} Ένα ράφι σε αμερικανικό κατάστημα έχει μήκος $4$ ft. Πόσες ίντσες είναι; Πόσα μέτρα είναι περίπου (με $1$ ft $\approx0{,}3048$ m);

\zoni{Το σκέφτομαι}{ίδια μονάδα, διπλασιάζω, αληθές ή ψευδές}

\ex Ο Βασίλης λέει ότι $0{,}8$ m $<75$ cm, επειδή το $0{,}8$ είναι μικρότερο από το $75$. Έχει δίκιο; Πώς θα του το εξηγούσες;

\ex Το $\mathrm{Μ}$ είναι το μέσο του $\mathrm{ΑΒ}=12$ cm και το $\mathrm{Ν}$ το μέσο του $\mathrm{ΑΜ}$. Πόσο είναι τα $\mathrm{ΑΝ}$ και $\mathrm{ΝΒ}$;

\ex Σωστό ή λάθος; α) Το χιλιοστόμετρο είναι μεγαλύτερη μονάδα από το εκατοστόμετρο. \quad β) Αν $\mathrm{Μ}$ μέσο του $\mathrm{ΑΒ}$ τότε $\mathrm{ΑΜ}=\mathrm{ΜΒ}$. \quad γ) Αν $\mathrm{ΜΑ}=\mathrm{ΜΒ}$ τότε το $\mathrm{Μ}$ είναι πάντα το μέσο του $\mathrm{ΑΒ}$. \quad δ) Η περίμετρος ενός σχήματος αλλάζει αν μετρήσω τις πλευρές με άλλη σειρά. \quad ε) $1$ m $=0{,}001$ km.

\ex Αν διπλασιάσουμε κάθε πλευρά ενός ορθογωνίου, πόσο γίνεται η περίμετρος; Θα ίσχυε το ίδιο και για ένα τρίγωνο; Δοκίμασε με παράδειγμα.

\zoni{Για το σπίτι}{μέτρησε και πρόσθεσε}

\ex Μέτρησε με τα βήματά σου και με μετροταινία το ίδιο μήκος στο σπίτι (π.χ. το διάδρομο) και υπολόγισε το μήκος του βήματός σου.

\ex Μέτρησε τις πλευρές του δωματίου σου και υπολόγισε την περίμετρο της κάτοψής του. Βρες και το μέσο ενός σχοινιού διπλώνοντάς το, και έλεγξε με μετροταινία.
